%
\section{\crowsFoot{C}{O1}{D}{M1}}%
%
\begin{frame}[t]%
\frametitle{\crowsFoot{C}{O1}{D}{M1}}%
\begin{itemize}%
\item Es gibt zwei Entitätstypen~C und~D.%
\item<2-> Jede Entität vom Typ~C muss mit genau einer Entität vom Typ~D verbunden sein.
\item<3-> Jede Entität vom Typ~D ist mit keiner oder einer Entität vom Typ~C verbunden.
\item<4-> Beispiel:\uncover<5->{%
\begin{itemize}%
%
\item Jede Entität vom Typ~\emph{Person} kann ein Mitarbeiter~(\emph{Faculty}) der Uni sein {\dots} oder auch nicht.\uncover<6->{ %
Jede Entität vom Typ~\emph{Faculty} muss mit genau einer Entität vom Typ~\emph{Person} verbunden sein.}%
\end{itemize}%
}%
\end{itemize}%
%
\locateGraphic{4-}{width=0.9\paperwidth}{graphics/CD}{0.05}{0.55}%
\end{frame}%
%
\begin{frame}%
\frametitle{Lösungen}%
\begin{itemize}%
\item Eine Lösung mit drei Tabellen ergibt hier gar keinen Sinn.%
%
\item<2-> Wir \emph{wissen}, dass jede Entität vom Typ~C mit genau einer Entität vom Typ~D verbunden sein \emph{muss}.%
%
\item<3-> Die Lösung ist also so ähnlich wie die zwei-Tabellen-Variante aus dem vorigen Abschnitt.%
\end{itemize}%
\end{frame}%
%
\begin{frame}%
\frametitle{Zwei-Tabellen Lösung}%
\parbox{0.435\linewidth}{\small%
\begin{itemize}%
%
\only<-4>{%
\item Wir brauchen eine Tabelle~\sqlil{c} für die Entitäten vom Typ~C.%
}%
%
\only<-5>{%
\item<2->  Wir brauchen auch eine Tabelle~\sqlil{d} für die Entitäten vom Typ~D.%
}%
%
\only<-6>{%
\item<3-> Beide Tabellen sind wie im vorigen Beispiel strukturiert, mit Ersatz-Primärschlüsseln~\sqlil{cid} und~\sqlil{did} sowie mit den Attributen~\sqlil{x} und~\sqlil{y}.%
}%
%
\item<4-> Weil jede Entität vom Typ~C mit genau einer Entität vom Typ~D verbunden seien muss, fügen wir die Spalte~\sqlil{fkdid} zur Tabelle~\sqlil{c} hinzu, um Fremdschlüssel von D-Entitäten zu speichern.%
%
\item<5-> Diese Spalte hat daher eine \sqlil{REFERENCES}-Einschränkung zum Fremdschlüsssel, die wir später mit \sqlil{ALTER TABLE} hinzufügen%
%
\item<6-> \alert{Anders} als vorher muss die Spalte~\sqlil{fkdid} auch \sqlil{NOT NULL} sein, denn jede Entität vom Typ~C \emph{muss} mit einer Entität vom Typ~D verbunden sein.%
%
\item<7-> Die Spalte hat auch eine \sqlil{UNIQUE}-Einschränkung, denn Entitäten vom Typ~D dürfen jeweils höchstens einen Eintrag aus der Tabelle~\sqlil{c} referenzieren.%
%
\end{itemize}%
}%
\gitLoadAndExecSQL{CD_tables}{}{conceptualToRelational}{CD_tables.sql}{relationships}{}{}%
\listingSQLandOutput{}{CD_tables}{}{0.47}{0.1}{0.529}{0.87}
\end{frame}%
%
\begin{frame}[t]%
\frametitle{Befüllen mit Daten}%
\parbox{0.435\linewidth}{\small%
\begin{itemize}%
\item Wir können nur Zeilen in die Tabelle~\sqlil{c} einfügen, die auf bereits existierende Zeilen in Tabelle~\sqlil{d} verweisen.
%
\item<2-> Die Konsequenz ist, dass diese Datensätze zuerst erzeugt werden müssen.%
%
\item<3-> Nach dem wir die Datensätze in die Tabelle~\sqlil{d} eingefügt haben, können wir Zeilen in die Tabelle~\sqlil{c} einfügen.%
%
\item<4-> Jede dieser Zeilen muss einen gültigen und einmaligen Fremdschlüssel~\sqlil{fkdid} haben, der einen Primärschlüsselwert~\sqlil{did} in Tabelle~\sqlil{d} referenziert.%
%
\item<5-> Die Daten können über einziges \sqlil{INNER JOIN} zusammengeführt werden.
\end{itemize}%
}%
\gitLoadAndExecSQL{CD_insert_and_select}{}{conceptualToRelational}{CD_insert_and_select.sql}{relationships}{}{}%
\listingSQLandOutput{}{CD_insert_and_select}{}{0.47}{0.1}{0.529}{0.87}
\end{frame}%
%
\begin{frame}[t]%
\frametitle{Der Inhalt der Zwei Tabellen}%
%
\gitExec{cdtrmTableC}{\databasesCodeRepo}{conceptualToRelational}{../_scripts_/db_table_to_latex_table.sh relationships c cid;fkdid;x}%
\gitExec{cdtrmTableD}{\databasesCodeRepo}{conceptualToRelational}{../_scripts_/db_table_to_latex_table.sh relationships d did;y}%
%
\vfill%
%
\begin{itemize}%
\item Dies sind die Daten in den zwei Tabellen.%
\end{itemize}%
\vfill%
%
\begin{center}%
\strut\hfill\strut%
\centering\input{\gitFile{cdtrmTableC}}%
\strut\hfill\strut%
\centering\input{\gitFile{cdtrmTableD}}%
\strut\hfill\strut%
\end{center}%
%
\hfill%
\end{frame}%
%
\begin{frame}[t]%
\frametitle{Testen der Einschränkungen}%
\parbox{0.435\linewidth}{\small%
\begin{itemize}%
%
\only<-4>{%
\item Es ist unmöglich, eine Zeile in Tabelle~\sqlil{c} zu haben, die mit mehr als einer Zeile in Tabelle~\sqlil{d} in Beziehung steht.%
\item<2-> Das Feld~\sqlil{fkdid} kann ja nur genau einen Wert haben.%
}%
%
\item<3-> Es ist auch unmöglich, eine Zeile in Tabelle~\sqlil{d} zu haben, die mit mehr als einer Zeile in Tabelle~\sqlil{c} in Beziehung steht.%
\item<4-> Die \sqlil{UNIQUE}-Einschränkung auf Spalte~\sqlil{fkdid} in Tabelle~\sqlil{c} verbietet das.%
%
\item<5-> Wir können auch keine Zeile in Tabelle~\sqlil{c} einfügen, die keine Zeile in Tabelle~\sqlil{d} referenziert.
%
\item<6-> Das wird durch die \sqlil{NOT NULL}-Einschränkung verhindert.%
%
\item<7-> Beachten Sie, dass daher Tabelle~\sqlil{d} niemals weniger Einträge als Tabelle~\sqlil{c} haben kann.%
%
\end{itemize}%
}%
%
\gitLoadAndExecSQL{CD_insert_error_1}{}{conceptualToRelational}{CD_insert_error_1.sql}{relationships}{}{}%
\gitLoadAndExecSQL{CD_insert_error_2}{}{conceptualToRelational}{CD_insert_error_2.sql}{relationships}{}{}%
\listingSQLandOutput{3-4}{CD_insert_error_1}{}{0.47}{0.1}{0.529}{0.87}
\listingSQLandOutput{5-}{CD_insert_error_2}{}{0.47}{0.1}{0.529}{0.87}
\end{frame}%
